\chapter{数值实验设定与结果对比分析}\label{chap:empirical} \% 确保已加载 tikz 包

本章基于第四章构建的 3/2 随机波动率模型九类期权定价方法，开展系统性数值实验，核心目标是量化对比不同方法在定价精度与计算效率两大核心维度的性能差异，验证不同方法在常规金融场景与极端参数下的适用性，为 3/2 模型的实际工程应用提供量化依据。所有实验均在统一的硬件环境下完成，结果可复现、可对比。

\section{数值实验环境与设定}

\subsection{硬件与软件环境}
本次数值实验的运行硬件为消费级笔记本联想电脑ThinkBook14，具体配置如下：处理器为 AMD Ryzen 7 6800H with Radeon Graphics，基础主频 3.20GHz，8 核 16 线程，标准 TDP 45W；内存为 16.0GB DDR5 双通道内存，实验运行时可用内存 13.7GB，可支撑大规模蒙特卡洛路径模拟与频域网格计算；操作系统为 Windows 11 企业版（版本号 26200.7840），64 位架构；开发与计算环境基于 Python 3.9 实现，核心数值计算依赖向量化运算库 numpy、特殊函数与概率分布库 scipy，所有方法均采用向量化实现以最大化利用硬件算力，未使用 GPU 加速，结果反映方法在通用 CPU 环境下的实际性能。

\subsection{测试案例设计}
本次实验共设计 7 组测试案例，覆盖了金融衍生品定价中最常见的常规场景与极端场景，包括平值（ATM）、实值（ITM）、虚值（OTM）期权，短期到期、高杠杆效应（强价格 - 波动率负相关）、高波动率的波动率（vol-of-vol）等典型场景，所有案例的基准价格均来自对应文献的精确解析结果或高精度数值解，作为精度评价的基准。7 组案例的具体参数与设计逻辑如下表所示：

\begin{table}[htbp ]
    \centering
    \caption{Test Cases and Parameter Configurations}\label{tab:cases}
    \resizebox{\textwidth}{!}{
    \begin{tabular}{lccccccccccc}
        \toprule
        Case & $\sigma$ & vov & $\rho$ & $\kappa$ & $\theta$ & $r$ & $T$ & $K$ & $S_0$ & 精确价格 ($p_\text{exact}$) & 场景描述 \\
        \midrule
        Case I & 1 & 0.2 & -0.5 & 2 & 1.5 & 0.05 & 1 & 1 & 1 & 0.443059 & Baldeaux (2012) \\
        \midrule
        Case II & 0.06 & 8.56 & -0.99 & 22.84 & 0.218 & 0 & 0.5 & 95 & 100 & 10.364 & Kouarfate et al. (2021), ATM \\
         &  &  &  &  &  &  &  & 100 & 100 & 7.386 & \\
         &  &  &  &  &  &  &  & 105 & 100 & 4.938 & \\
        \midrule
        Case III & 0.06 & 3.2 & -0.99 & 20.48 & 0.218 & 0 & 0.5 & 95 & 100 & 11.724 & Kouarfate et al. (2021), ATM \\
         &  &  &  &  &  &  &  & 100 & 100 & 8.999 & \\
         &  &  &  &  &  &  &  & 105 & 100 & 6.710 & \\
        \midrule
        Case IV & 1 & 0.3 & 0.5 & 0.7 & 0.6 & 0 & 0.1 & 47 & 49 & 7.040 & Near expiration and ATM \\
         &  &  &  &  &  &  &  & 48 & 49 & 6.500 & \\
         &  &  &  &  &  &  &  & 49 & 49 & 6.121 & \\
         &  &  &  &  &  &  &  & 50 & 49 & 5.7006 & \\
         &  &  &  &  &  &  &  & 51 & 49 & 5.305 & \\
        \midrule
        Case V & 1 & 0.3 & 0.5 & 0.7 & 0.6 & 0 & 0.1 & 10 & 50 & 39.9985 & Near exp only \\
         &  &  &  &  &  &  &  & 20 & 50 & 30.002 & \\
         &  &  &  &  &  &  &  & 40 & 50 & 11.9266 & \\
        \midrule
        Case VI & 1 & 3.3 & 0.5 & 0.7 & 0.6 & 0 & 1 & 47 & 49 & 12.7468 & ATM only \\
         &  &  &  &  &  &  &  & 48 & 49 & 12.3995 & \\
         &  &  &  &  &  &  &  & 49 & 49 & 12.0658 & \\
         &  &  &  &  &  &  &  & 50 & 49 & 11.7452 & \\
         &  &  &  &  &  &  &  & 51 & 49 & 11.4371 & \\
        \midrule
        Case VII & 0.06 & 3.2 & -0.99 & 20.48 & 0.218 & 0 & 0.5 & 95 & 100 & 11.7235 & Lewis (2000), ATM \\
         &  &  &  &  &  &  &  & 100 & 100 & 8.9978 & \\
         &  &  &  &  &  &  &  & 105 & 100 & 6.7091 & \\
        \bottomrule
    \end{tabular}
    }
\end{table}
Case I 为\citet{baldeauxEXACTSIMULATION322012} 中的经典基准算例，用于验证方法的基础理论定价精度；Case II 为极端高 vol-of-vol、强价格 - 波动率负相关场景，用于模拟市场危机下的高杠杆定价环境；Case III 为中高 vol-of-vol、强负相关场景，与 Case II 形成对照，检验 vol-of-vol 参数对方法性能的影响；Case IV 为短期到期、近平值场景，对应高频交易中的短期期权定价需求；Case V 为短期到期、深度实值场景，用于检验方法在深度实值期权下的定价稳定性；Case VI 为高 vol-of-vol、平值场景，用于检验方法在中等期限、高波动环境下的长期定价性能；Case VII 为 Lewis (2000) 提出的经典算例，作为频域定价方法的基准验证场景。这 7 组案例完整覆盖了 3/2 模型在实际应用中的绝大多数场景，既包含了用于验证基础精度的基准算例，也包含了用于测试方法鲁棒性的极端参数场景，可全面对比不同方法的性能边界。
\subsection{评价指标定义}
本次实验构建三级评价指标体系，以量化评估方法的综合性能，指标优先级从高到低依次为数值稳定性、定价精度、计算效率。数值稳定性为方法具备工程实用性的基础判定标准，若方法输出结果为非数值型（nan）或无穷大（inf），判定为该场景下数值失效，丧失定价能力。定价精度采用绝对偏差与相对偏差两个指标衡量，绝对偏差为方法计算价格与基准价格的差值，相对偏差为绝对偏差与基准价格的百分比，用于不同执行价、不同价格水平下的精度横向对比。计算效率采用单组案例（含对应全部执行价）的完整运行时长衡量，单位为秒，反映方法的算力需求与工程应用适配性。

\section{时间步进类条件蒙特卡洛方法结果分析}

时间步进类方法共包含 4 种方案：Euler/Milstein 离散格式、精确步进（Exact Stepping）、Poisson-Gamma 近似精确步进、QE 近似精确步进，实验设置了 3 组时间步长 dt=1/50、1/500、1/5000，对应从粗到细的时间离散精度，验证步长对方法性能的影响，完整结果如下表所示：

\begin{table}[htbp]
  \centering
  \caption{Time-Stepping Conditional Monte Carlo Results}\label{tab:timestepping}
  \resizebox{\textwidth}{!}{%
  \begin{tabular}{lclclcccccccc}
    \toprule
    \multirow{2}{*}{Case} & \multirow{2}{*}{$dt$} & \multirow{2}{*}{Strike} & \multicolumn{2}{c}{Euler/Milstein scheme} & \multicolumn{2}{c}{Exact Stepping} & \multicolumn{2}{c}{Poisson-Gamma Almost Exact Stepping} & \multicolumn{2}{c}{QE Almost Exact Stepping} \\
    \cmidrule(lr){4-5} \cmidrule(lr){6-7} \cmidrule(lr){8-9} \cmidrule(lr){10-11}
    & & & Time (s) & Option Bias & Time (s) & Option Bias & Time (s) & Option Bias & Time (s) & Option Bias \\
    \midrule
    \multirow{3}{*}{Case I}
    & 1/50   & 1 & 0.405223    & $-0.01165$ ($-2.63\%$)     & 0.428566     & $-0.0009242$ ($-0.21\%$)  & 0.6209759   & $-0.0005159$ ($-0.12\%$)  & 0.6758925   & $-0.0002905$ ($-0.07\%$) \\
    & 1/500  & 1 & 2.8924741   & $-0.001347$ ($-0.30\%$)    & 3.848103     & $-0.000335$ ($-0.07\%$)   & 5.5486124   & $-0.0004362$ ($-0.10\%$)  & 6.66017     & $-0.0001411$ ($-0.03\%$) \\
    & 1/5000 & 1 & 29.065831   & $0.0002249$ ($0.06\%$)     & 40.3844027   & $-0.0001585$ ($-0.04\%$)  & 88.8646271  & $-0.000115$ ($-0.02\%$)    & 101.9920107 & $-0.000367$ ($-0.04\%$) \\
    \midrule
    \multirow{9}{*}{Case II}
    & \multirow{3}{*}{1/50}  & 95  & \multirow{3}{*}{0.2677657} & $230.2249$ ($-9.06\%$)   & \multirow{3}{*}{0.2455173} & $-0.0114$ ($-0.11\%$)     & \multirow{3}{*}{0.4654895} & $-0.004389$ ($-0.05\%$)   & \multirow{3}{*}{0.390341} & $-0.000977$ ($-0.01\%$) \\
    &                         & 100 &                           & $2.302$ ($31.17\%$)      &                           & $-0.0079$ ($-0.23\%$)     &                           & $-0.007585$ ($-0.10\%$)   &                           & $-0.01131$ ($-0.15\%$) \\
    &                         & 105 &                           & $7.401$ ($714.93\%$)     &                           & $0.001158$ ($0.02\%$)     &                           & $0.002399$ ($0.05\%$)     &                           & $0.01003$ ($0.20\%$) \\
    \cmidrule(lr){2-11}
    & \multirow{3}{*}{1/500} & 105 & \multirow{3}{*}{1.504743}  & $7.427$ ($150.41\%$)    & \multirow{3}{*}{1.9639806} & $0.0002409$ ($0.00\%$)    & \multirow{3}{*}{2.9796648} & $-0.004103$ ($-0.08\%$)   & \multirow{3}{*}{3.397378} & $0.01213$ ($0.25\%$) \\
    &                         & 100 &                           & $0.15086$ ($3.89\%$)     &                           & $0.003497$ ($0.03\%$)     &                           & $-0.01347$ ($-0.13\%$)    &                           & $0.01355$ ($0.18\%$) \\
    &                         & 95  &                           & $0.0909$ ($1.23\%$)      &                           & $-0.001302$ ($-0.02\%$)   &                           & $-0.01213$ ($-0.16\%$)    &                           & $-0.01645$ ($-0.22\%$) \\
    \cmidrule(lr){2-11}
    & \multirow{3}{*}{1/5000}& 105 & \multirow{3}{*}{14.525341} & $0.0909$ ($1.89\%$)     & \multirow{3}{*}{20.2330929}& $-0.002073$ ($-0.06\%$)   & \multirow{3}{*}{29.5629414}& $-0.01034$ ($-0.21\%$)    & \multirow{3}{*}{48.5850875}& $-0.01676$ ($-0.34\%$) \\
    &                         & 100 &                           & $0.36318$ ($1.23\%$)     &                           & $-0.00273$ ($-0.07\%$)    &                           & $-0.007066$ ($-0.10\%$)   &                           & $-0.0309$ ($-0.46\%$) \\
    &                         & 95  &                           & $0.093112$ ($0.89\%$)    &                           & $-0.00142$ ($-0.01\%$)    &                           & $-0.009375$ ($-0.34\%$)   &                           & $-0.02017$ ($-0.17\%$) \\
    \midrule
    \multirow{9}{*}{Case III}
    & \multirow{3}{*}{1/50}  & 95  & \multirow{3}{*}{0.238946}  & $364.9$ ($9054.87\%$)   & \multirow{3}{*}{0.2417619} & $-0.0126$ ($-0.14\%$)     & \multirow{3}{*}{0.3870248} & $-0.05426$ ($-0.66\%$)    & \multirow{3}{*}{0.4423578} & $-0.03312$ ($-0.37\%$) \\
    &                         & 100 &                           & $36.4$ ($5438.78\%$)     &                           & $-0.005401$ ($-0.08\%$)   &                           & $-0.0472$ ($-0.70\%$)     &                           & $-0.0326$ ($-0.50\%$) \\
    &                         & 105 &                           & $0.05495$ ($-0.57\%$)    &                           & $-0.0107$ ($-0.12\%$)     &                           & $-0.0188$ ($-0.21\%$)     &                           & $0.01358$ ($0.17\%$) \\
    \cmidrule(lr){2-11}
    & \multirow{3}{*}{1/500} & 105 & \multirow{3}{*}{1.4755198} & $-0.05096$ ($-0.72\%$)  & \multirow{3}{*}{1.9966579} & $-0.009423$ ($-0.14\%$)   & \multirow{3}{*}{2.8634368} & $-0.01699$ ($-0.25\%$)    & \multirow{3}{*}{3.3736463} & $0.0144$ ($0.20\%$) \\
    &                         & 100 &                           & $-0.02709$ ($-0.23\%$)   &                           & $-0.00122$ ($-0.01\%$)    &                           & $-0.00975$ ($-0.11\%$)    &                           & $-0.0155$ ($-0.17\%$) \\
    &                         & 95  &                           & $-0.02338$ ($-0.33\%$)   &                           & $-0.006346$ ($-0.07\%$)   &                           & $-0.0375$ ($-0.34\%$)     &                           & $-0.02017$ ($-0.17\%$) \\
    \cmidrule(lr){2-11}
    & \multirow{3}{*}{1/5000}& 105 & \multirow{3}{*}{14.4603692}& $-0.02866$ ($-0.40\%$)  & \multirow{3}{*}{19.967838} & $-0.008676$ ($-0.13\%$)   & \multirow{3}{*}{28.5484009}& $-0.00941$ ($-0.14\%$)    & \multirow{3}{*}{48.5172484}& $-0.02275$ ($-0.31\%$) \\
    &                         & 100 &                           & $-0.01184$ ($-0.17\%$)   &                           & $-0.01091$ ($-0.16\%$)    &                           & $-0.03466$ ($-0.52\%$)    &                           & $-0.0207$ ($-0.25\%$) \\
    &                         & 95  &                           & $0.018024$ ($0.24\%$)    &                           & $-0.005735$ ($-0.06\%$)   &                           & $-0.00736$ ($0.11\%$)     &                           & $-0.00109$ ($-0.02\%$) \\
    \midrule
    \multirow{15}{*}{Case IV}
    & \multirow{5}{*}{1/50}  & 48 & \multirow{5}{*}{0.0891715} & $0.01468$ ($0.26\%$)     & \multirow{5}{*}{0.0836116} & $-0.008412$ ($-0.15\%$)   & \multirow{5}{*}{0.1061546} & $0.004264$ ($0.07\%$)     & \multirow{5}{*}{0.1498051} & $-0.0002527$ ($-0.00\%$) \\
    &                        & 49 &                          & $0.01468$ ($0.23\%$)     &                          & $-0.007911$ ($-0.13\%$)   &                          & $0.004212$ ($0.07\%$)     &                          & $-0.0002652$ ($-0.00\%$) \\
    &                        & 50 &                          & $0.01468$ ($0.26\%$)     &                          & $0.006557$ ($0.10\%$)     &                          & $0.003103$ ($0.04\%$)     &                          & $-0.004188$ ($-0.06\%$) \\
    &                        & 51 &                          & $0.00985$ ($-0.04\%$)    &                          & $0.007421$ ($0.14\%$)     &                          & $0.00358$ ($0.09\%$)      &                          & $-0.00351$ ($-0.08\%$) \\
    &                        & 47 &                          & $-0.00295$ ($-0.04\%$)   &                          & $-0.01091$ ($-0.16\%$)    &                          & $0.007103$ ($0.10\%$)     &                          & $-0.003395$ ($-0.05\%$) \\
    \cmidrule(lr){2-11}
    & \multirow{5}{*}{1/500} & 49 & \multirow{5}{*}{0.323018}  & $0.006214$ ($0.10\%$)    & \multirow{5}{*}{0.4360611} & $0.00662$ ($0.11\%$)      & \multirow{5}{*}{0.6052143} & $0.003458$ ($0.06\%$)     & \multirow{5}{*}{0.7028849} & $0.00635$ ($0.09\%$) \\
    &                        & 48 &                          & $-0.001867$ ($-0.03\%$)  &                          & $0.006706$ ($0.12\%$)     &                          & $0.003327$ ($0.06\%$)     &                          & $-0.003306$ ($-0.06\%$) \\
    &                        & 50 &                          & $-0.001804$ ($-0.03\%$)  &                          & $0.006374$ ($0.12\%$)     &                          & $0.00312$ ($0.06\%$)      &                          & $-0.003293$ ($-0.06\%$) \\
    &                        & 51 &                          & $-0.001804$ ($-0.03\%$)  &                          & $0.006539$ ($0.12\%$)     &                          & $0.002685$ ($0.04\%$)     &                          & $-0.003207$ ($-0.06\%$) \\
    &                        & 47 &                          & $0.006706$ ($0.10\%$)    &                          & $0.006259$ ($0.09\%$)     &                          & $0.00703$ ($0.10\%$)      &                          & $0.00693$ ($0.10\%$) \\
    \cmidrule(lr){2-11}
    & \multirow{5}{*}{1/5000}& 49 & \multirow{5}{*}{3.0013151} & $0.00152$ ($0.03\%$)     & \multirow{5}{*}{4.0081004} & $-0.004082$ ($-0.06\%$)   & \multirow{5}{*}{5.7636176} & $0.00318$ ($0.05\%$)      & \multirow{5}{*}{10.02923375} & $0.001356$ ($0.00\%$) \\
    &                         & 48 &                          & $0.00152$ ($0.03\%$)     &                          & $-0.00989$ ($-0.96\%$)    &                          & $0.003308$ ($0.05\%$)     &                          & $0.000693$ ($0.01\%$) \\
    &                         & 50 &                          & $0.0002269$ ($0.00\%$)   &                          & $-0.003774$ ($-0.07\%$)   &                          & $0.003322$ ($0.06\%$)     &                          & $8.264\times 10^{-5}$ ($0.00\%$) \\
    &                         & 51 &                          & $-0.0002143$ ($-0.00\%$) &                          & $-0.00356$ ($-0.07\%$)    &                          & $0.003245$ ($0.06\%$)     &                          & $6.51\times 10^{-5}$ ($0.00\%$) \\
    &                         & 47 &                          & $-0.001015$ ($-0.01\%$)  &                          & $-0.01999$ ($-0.30\%$)    &                          & $0.005225$ ($0.07\%$)     &                          & $0.0005575$ ($0.01\%$) \\
    \midrule
    \multirow{9}{*}{Case V}
    & \multirow{3}{*}{1/50}  & 20 & \multirow{3}{*}{0.0873017} & $-0.00994$ ($-0.03\%$)   & \multirow{3}{*}{0.0845473} & $-0.01497$ ($-0.07\%$)    & \multirow{3}{*}{0.1103219} & $0.005325$ ($0.02\%$)     & \multirow{3}{*}{0.129903} & $0.000554$ ($0.00\%$) \\
    &                         & 40 &                           & $0.005247$ ($0.04\%$)     &                           & $0.01425$ ($0.12\%$)      &                           & $0.00362$ ($0.02\%$)      &                           & $0.000675$ ($0.01\%$) \\
    &                         & 10 &                           & $-0.002526$ ($-0.01\%$)  &                           & $0.01425$ ($0.04\%$)      &                           & $0.008478$ ($0.02\%$)     &                           & $-0.00144$ ($-0.00\%$) \\
    \cmidrule(lr){2-11}
    & \multirow{3}{*}{1/500} & 40 & \multirow{3}{*}{0.3298876} & $-0.001768$ ($-0.01\%$)  & \multirow{3}{*}{0.4290495} & $0.01129$ ($0.09\%$)      & \multirow{3}{*}{0.5894391} & $0.00529$ ($0.03\%$)      & \multirow{3}{*}{0.7007175} & $-0.001399$ ($-0.02\%$) \\
    &                         & 20 &                           & $0.002483$ ($0.01\%$)     &                           & $0.008224$ ($0.05\%$)     &                           & $0.006829$ ($0.06\%$)     &                           & $-0.00217$ ($-0.02\%$) \\
    &                         & 10 &                           & $0.00053$ ($0.00\%$)      &                           & $-0.008224$ ($-0.02\%$)   &                           & $0.007159$ ($0.02\%$)     &                           & $0.000674$ ($0.00\%$) \\
    \cmidrule(lr){2-11}
    & \multirow{3}{*}{1/5000}& 40 & \multirow{3}{*}{2.9262788} & $0.0005752$ ($0.00\%$)   & \multirow{3}{*}{4.0733267} & $-0.008224$ ($-0.02\%$)   & \multirow{3}{*}{5.7201683} & $0.007159$ ($0.02\%$)     & \multirow{3}{*}{9.9833601} & $0.0008315$ ($0.01\%$) \\
    &                         & 20 &                           & $0.0005752$ ($0.00\%$)   &                           & $-0.008224$ ($-0.03\%$)   &                           & $0.007159$ ($0.02\%$)     &                           & $0.000687$ ($0.00\%$) \\
    &                         & 10 &                           & $0.000949$ ($0.00\%$)     &                           & $-0.005735$ ($-0.01\%$)   &                           & $0.00539$ ($0.01\%$)      &                           & $0.000933$ ($0.00\%$) \\
    \midrule
    \multirow{15}{*}{Case VI}
    & \multirow{5}{*}{1/50}  & 48 & \multirow{5}{*}{0.3831229} & $\inf$ ($\inf\%$)    & \multirow{5}{*}{0.4302237} & $225.7$ ($171.53\%$)      & \multirow{5}{*}{0.7873058} & $36.17$ ($283.9\%$)       & \multirow{5}{*}{0.7045123} & $\inf$ ($\inf\%$) \\
    &                        & 49 &                          & $\inf$ ($\inf\%$)    &                          & $225.6$ ($1871.65\%$)     &                          & $36.17$ ($300.01\%$)      &                          & $\inf$ ($\inf\%$) \\
    &                        & 50 &                          & $\inf$ ($\inf\%$)    &                          & $225.6$ ($1922.8\%$)      &                          & $36.17$ ($308.2\%$)       &                          & $\inf$ ($\inf\%$) \\
    &                        & 51 &                          & $\inf$ ($\inf\%$)    &                          & $225.6$ ($1974.65\%$)     &                          & $36.17$ ($316.52\%$)      &                          & $\inf$ ($\inf\%$) \\
    &                        & 47 &                          & $\inf$ ($\inf\%$)    &                          & $225.7$ ($1771.21\%$)     &                          & $36.17$ ($291.99\%$)      &                          & $\inf$ ($\inf\%$) \\
    \cmidrule(lr){2-11}
    & \multirow{5}{*}{1/500} & 49 & \multirow{5}{*}{2.8947977} & $\inf$ ($\inf\%$)    & \multirow{5}{*}{3.8898247} & $-0.0573$ ($-0.46\%$)     & \multirow{5}{*}{5.8676512} & $0.01475$ ($0.12\%$)      & \multirow{5}{*}{6.6754493} & $4.106\times 10^{68}$ (nan\%) \\
    &                        & 48 &                          & $\inf$ ($\inf\%$)    &                          & $-0.0575$ ($-0.47\%$)     &                          & $0.01395$ ($0.11\%$)      &                          & $4.106\times 10^{68}$ (nan\%) \\
    &                        & 50 &                          & $\inf$ ($\inf\%$)    &                          & $-0.05727$ ($-0.49\%$)    &                          & $0.01336$ ($0.11\%$)      &                          & $4.106\times 10^{68}$ (nan\%) \\
    &                        & 51 &                          & $\inf$ ($\inf\%$)    &                          & $-0.05716$ ($-0.50\%$)    &                          & $0.01297$ ($0.11\%$)      &                          & $4.106\times 10^{68}$ (nan\%) \\
    &                        & 47 &                          & $\inf$ ($\inf\%$)    &                          & $-0.05716$ ($-0.47\%$)    &                          & $0.0136$ ($0.11\%$)       &                          & $4.106\times 10^{68}$ (nan\%) \\
    \cmidrule(lr){2-11}
    & \multirow{5}{*}{1/5000}& 49 & \multirow{5}{*}{29.229539} & $468.6$ ($3678.98\%$)   & \multirow{5}{*}{55.4278611} & $0.00935$ ($0.08\%$)      & \multirow{5}{*}{57.7701995} & $0.02743$ ($0.22\%$)      & \multirow{5}{*}{99.3793894} & $3.057\times 10^{21}$ (nan\%) \\
    &                         & 48 &                          & $468.6$ ($3782.16\%$)   &                          & $0.00958$ ($0.08\%$)      &                          & $0.02739$ ($0.22\%$)      &                          & $3.057\times 10^{21}$ (nan\%) \\
    &                         & 50 &                          & $468.6$ ($3893.10\%$)   &                          & $0.00962$ ($0.08\%$)      &                          & $0.02713$ ($0.23\%$)      &                          & $3.057\times 10^{21}$ (nan\%) \\
    &                         & 51 &                          & $468.6$ ($4100.83\%$)   &                          & $0.00943$ ($0.08\%$)      &                          & $0.02707$ ($0.24\%$)      &                          & $3.057\times 10^{21}$ (nan\%) \\
    &                         & 47 &                          & $468.6$ ($3678.98\%$)   &                          & $0.00958$ ($0.08\%$)      &                          & $0.02743$ ($0.22\%$)      &                          & $3.057\times 10^{21}$ (nan\%) \\
    \midrule
    \multirow{9}{*}{Case VII}
    & \multirow{3}{*}{1/50}  & 95  & \multirow{3}{*}{0.2103698} & $364.9$ ($4055.43\%$)   & \multirow{3}{*}{0.2497061} & $-0.0114$ ($-0.13\%$)     & \multirow{3}{*}{0.3766218} & $-0.05822$ ($-0.66\%$)    & \multirow{3}{*}{0.3742216} & $-0.03192$ ($-0.36\%$) \\
    &                         & 100 &                           & $36.4$ ($5439.54\%$)     &                           & $-0.004501$ ($-0.07\%$)   &                           & $-0.0436$ ($-0.65\%$)     &                           & $-0.03236$ ($-0.48\%$) \\
    &                         & 105 &                           & $0.05495$ ($-0.46\%$)    &                           & $-0.001013$ ($-0.01\%$)   &                           & $-0.0198$ ($-0.16\%$)     &                           & $0.015$ ($0.13\%$) \\
    \cmidrule(lr){2-11}
    & \multirow{3}{*}{1/500} & 105 & \multirow{3}{*}{1.4952733} & $-0.04372$ ($-0.71\%$)  & \multirow{3}{*}{1.9388048} & $-0.009498$ ($-0.11\%$)   & \multirow{3}{*}{2.9082565} & $-0.0176$ ($-0.20\%$)     & \multirow{3}{*}{3.36681} & $0.01678$ ($0.21\%$) \\
    &                         & 100 &                           & $-0.01256$ ($-0.18\%$)   &                           & $-0.01523$ ($-0.13\%$)    &                           & $-0.01609$ ($-0.24\%$)    &                           & $0.01431$ ($0.19\%$) \\
    &                         & 95  &                           & $-0.02709$ ($-0.23\%$)   &                           & $-0.00742$ ($-0.01\%$)    &                           & $-0.03935$ ($-0.42\%$)    &                           & $-0.01967$ ($-0.17\%$) \\
    \cmidrule(lr){2-11}
    & \multirow{3}{*}{1/5000}& 105 & \multirow{3}{*}{14.7206092}& $-0.02589$ ($-0.31\%$)  & \multirow{3}{*}{34.2497818}& $-0.005146$ ($-0.06\%$)   & \multirow{3}{*}{48.1214847}& $-0.0321$ ($-0.50\%$)     & \multirow{3}{*}{52.2054604}& $-0.02155$ ($-0.24\%$) \\
    &                         & 100 &                           & $-0.0218$ ($-0.39\%$)    &                           & $-0.00076$ ($-0.01\%$)    &                           & $-0.03376$ ($-0.40\%$)    &                           & $-0.01987$ ($-0.22\%$) \\
    &                         & 95  &                           & $-0.02596$ ($-0.23\%$)   &                           & $-0.009742$ ($-0.12\%$)   &                           & $-0.03621$ ($-0.34\%$)    &                           & $-0.01931$ ($-0.17\%$) \\
    \bottomrule
  \end{tabular}%
  }
\end{table}
实验结果显示，四类方法在数值稳定性与定价性能上呈现显著分化。Euler/Milstein 离散格式的核心局限在于对时间步长的高度敏感性与极端场景下的数值失效。在 1/50 粗步长下，该方法在 Case II 虚值期权场景中相对偏差达 714.93\%，在 Case VI 高 vol-of-vol 场景中全部输出无穷大结果，完全丧失定价能力；当步长缩小至 1/500 时，该方法在 Case VI 场景中仍全部数值失效，Case II 平值期权的相对偏差仍达 3.89\%，无法满足工程定价的精度要求；仅当步长进一步缩小至 1/5000 时，该方法在常规场景下的相对偏差可降至 1\% 以内，但对应计算时间从 0.2 秒量级飙升至 30 秒以上，时间成本提升两个数量级，计算效率极低。该现象的核心成因在于 3/2 模型方差过程的强非线性特征，粗步长下 Euler 格式的一阶离散误差被显著放大，易出现负方差的非物理结果，最终导致定价失效。
与 Euler/Milstein 格式不同的是，精确步进、Poisson-Gamma 近似精确步进、QE 近似精确步进三类方法，在全部 7 组测试案例、所有执行价与所有时间步长下，均未出现数值失效，始终可输出有效定价结果，是本次实验中唯一具备全场景稳定性的蒙特卡洛类方法。即使在 1/50 粗步长下，精确步进方法在 Case II 平值期权场景中的相对偏差仅为 - 0.03\%，Poisson-Gamma 与 QE 方法的偏差也控制在 0.2\% 以内，精度显著优于 Euler 格式 1/5000 细步长下的结果，而计算时间仅为 0.2 秒左右，不足 Euler 细步长计算时间的 1\%。在 Case VI 高 vol-of-vol 极端场景中，1/500 步长下精确步进方法的相对偏差仅为 - 0.46\%，Poisson-Gamma 方法仅为 0.12\%，而同步长下 Euler 格式完全数值失效。三类方法的性能差异同样可通过实验结果量化验证，精确步进方法的定价精度最高，所有场景下的相对偏差均控制在 0.5\% 以内，无时间离散偏差，为时间步进类方法的精度基准；Poisson-Gamma 与 QE 方法的精度略低于精确步进方法，但计算时间与 Euler 格式处于同一量级，在极端场景下仍保持数值稳定性，实现了定价精度与计算效率的最优平衡。

\section{终端精确与近似精确模拟方法结果分析}
终端模拟类方法完全规避了时间离散过程，直接针对到期方差与积分方差的联合分布抽样，共包含 5 种方案：Baldeaux精确模拟、基于 Fourier-Cosine 的精确模拟、基于矩匹配的精确模拟、逆高斯（IG）近似精确模拟、Gamma 近似精确模拟，完整结果如下表所示：

\begin{table}[htbp]
  \centering
  \caption{Terminal Exact and Almost Exact Simulation Results}
  \label{tab:terminal}
  \resizebox{\textwidth}{!}{%
  \begin{tabular}{lclcccccccccc}
    \toprule
    \multirow{2}{*}{Case} & \multirow{2}{*}{Strike} & \multicolumn{2}{c}{Exact Simulation} & \multicolumn{2}{c}{Exact Simulation (Use Fourier-Cosine)} & \multicolumn{2}{c}{Exact Simulation (Use Moment Matching)} & \multicolumn{2}{c}{Almost Exact Simulation (Inverse Gaussian)} & \multicolumn{2}{c}{Almost Exact Simulation (Gamma)} \\
    \cmidrule(lr){3-4} \cmidrule(lr){5-6} \cmidrule(lr){7-8} \cmidrule(lr){9-10} \cmidrule(lr){11-12}
    & & Time(s) & Option Bias & Time(s) & Option Bias & Time(s) & Option Bias & Time(s) & Option Bias & Time(s) & Option Bias \\
    \midrule
    Case I & 1 & 59.5239847 & $-1.571\times10^{-5}\ (-0.00\%)$ & 25.4983391 & $-0.04977\ (-11.23\%)$ & 0.025079 & $-0.05216\ (-11.77\%)$ & 2.6390925 & $0.001326\ (0.30\%)$ & 2.164342 & $0.0008747\ (0.20\%)$ \\
    Case I & 95 & & $\mathrm{NaN}\ (\mathrm{NaN}\%)$ & & $\mathrm{NaN}\ (\mathrm{NaN}\%)$ & & $\mathrm{NaN}\ (\mathrm{NaN}\%)$ & & $-8.215\ (-79.26\%)$ & & $-7.956\ (-76.77\%)$ \\
    \midrule
    \multirow{3}{*}{Case II} & 100 & \multirow{3}{*}{69.0326374} & $\mathrm{NaN}\ (\mathrm{NaN}\%)$ & \multirow{3}{*}{8.3304516} & $\mathrm{NaN}\ (\mathrm{NaN}\%)$ & \multirow{3}{*}{0.034264} & $\mathrm{NaN}\ (\mathrm{NaN}\%)$ & \multirow{3}{*}{0.2152987} & $-6.213\ (-84.12\%)$ & \multirow{3}{*}{0.204456} & $-5.994\ (-81.16\%)$ \\
    & 105 & & $\mathrm{NaN}\ (\mathrm{NaN}\%)$ & & $\mathrm{NaN}\ (\mathrm{NaN}\%)$ & & $\mathrm{NaN}\ (\mathrm{NaN}\%)$ & & $-4.361\ (-88.31\%)$ & & $-4.195\ (-84.95\%)$ \\
    & 95 & & $\mathrm{NaN}\ (\mathrm{NaN}\%)$ & & $\mathrm{NaN}\ (\mathrm{NaN}\%)$ & & $\mathrm{NaN}\ (\mathrm{NaN}\%)$ & & $-11.64\ (-89.30\%)$ & & $-11.97\ (-99.56\%)$ \\
    \midrule
    \multirow{3}{*}{Case III} & 100 & \multirow{3}{*}{73.7386331} & $\mathrm{NaN}\ (\mathrm{NaN}\%)$ & \multirow{3}{*}{6.6512523} & $\mathrm{NaN}\ (\mathrm{NaN}\%)$ & \multirow{3}{*}{0.032648} & $\mathrm{NaN}\ (\mathrm{NaN}\%)$ & \multirow{3}{*}{0.256394} & $-8.96\ (-99.57\%)$ & \multirow{3}{*}{0.241701} & $-8.975\ (-99.73\%)$ \\
    & 105 & & $\mathrm{NaN}\ (\mathrm{NaN}\%)$ & & $\mathrm{NaN}\ (\mathrm{NaN}\%)$ & & $\mathrm{NaN}\ (\mathrm{NaN}\%)$ & & $-6.693\ (-99.74\%)$ & & $-6.699\ (-99.84\%)$ \\
    & 95 & & $\mathrm{NaN}\ (\mathrm{NaN}\%)$ & & $\mathrm{NaN}\ (\mathrm{NaN}\%)$ & & $\mathrm{NaN}\ (\mathrm{NaN}\%)$ & & $-11.64\ (-99.30\%)$ & & $-11.975\ (-99.56\%)$ \\
    \midrule
    \multirow{5}{*}{Case IV} & 47 & \multirow{5}{*}{79.0467981} & $\mathrm{NaN}\ (\mathrm{NaN}\%)$ & \multirow{5}{*}{6.5286679} & $\mathrm{NaN}\ (\mathrm{NaN}\%)$ & \multirow{5}{*}{0.026266} & $\mathrm{NaN}\ (\mathrm{NaN}\%)$ & \multirow{5}{*}{20.5061735} & $-5.559\ (-78.96\%)$ & \multirow{5}{*}{13.86428} & $38.31\ (544.13\%)$ \\
    & 48 & & $\mathrm{NaN}\ (\mathrm{NaN}\%)$ & & $\mathrm{NaN}\ (\mathrm{NaN}\%)$ & & $\mathrm{NaN}\ (\mathrm{NaN}\%)$ & & $-5.287\ (-81.34\%)$ & & $37.91\ (583.29\%)$ \\
    & 49 & & $\mathrm{NaN}\ (\mathrm{NaN}\%)$ & & $\mathrm{NaN}\ (\mathrm{NaN}\%)$ & & $\mathrm{NaN}\ (\mathrm{NaN}\%)$ & & $-5.135\ (-83.88\%)$ & & $37.37\ (610.49\%)$ \\
    & 50 & & $\mathrm{NaN}\ (\mathrm{NaN}\%)$ & & $\mathrm{NaN}\ (\mathrm{NaN}\%)$ & & $\mathrm{NaN}\ (\mathrm{NaN}\%)$ & & $-4.903\ (-86.01\%)$ & & $36.87\ (646.82\%)$ \\
    & 51 & & $\mathrm{NaN}\ (\mathrm{NaN}\%)$ & & $\mathrm{NaN}\ (\mathrm{NaN}\%)$ & & $\mathrm{NaN}\ (\mathrm{NaN}\%)$ & & $-4.664\ (-87.91\%)$ & & $36.36\ (685.41\%)$ \\
    \midrule
    \multirow{3}{*}{Case V} & 10 & \multirow{3}{*}{75.2066147} & $\mathrm{NaN}\ (\mathrm{NaN}\%)$ & \multirow{3}{*}{6.0646468} & $\mathrm{NaN}\ (\mathrm{NaN}\%)$ & \multirow{3}{*}{0.025757} & $\mathrm{NaN}\ (\mathrm{NaN}\%)$ & \multirow{3}{*}{22.739208} & $-5.255\ (-13.14\%)$ & \multirow{3}{*}{13.84294} & $43.83\ (109.58\%)$ \\
    & 20 & & $\mathrm{NaN}\ (\mathrm{NaN}\%)$ & & $\mathrm{NaN}\ (\mathrm{NaN}\%)$ & & $\mathrm{NaN}\ (\mathrm{NaN}\%)$ & & $-5.259\ (-17.53\%)$ & & $43.83\ (146.08\%)$ \\
    & 40 & & $\mathrm{NaN}\ (\mathrm{NaN}\%)$ & & $\mathrm{NaN}\ (\mathrm{NaN}\%)$ & & $\mathrm{NaN}\ (\mathrm{NaN}\%)$ & & $-6.356\ (-53.29\%)$ & & $41.99\ (352.08\%)$ \\
    \midrule
    \multirow{5}{*}{Case VI} & 47 & \multirow{5}{*}{69.5002801} & $\mathrm{NaN}\ (\mathrm{NaN}\%)$ & \multirow{5}{*}{21.3088199} & $\mathrm{NaN}\ (\mathrm{NaN}\%)$ & \multirow{5}{*}{0.033759} & $\mathrm{NaN}\ (\mathrm{NaN}\%)$ & \multirow{5}{*}{0.1932632} & $-0.06061\ (-0.48\%)$ & \multirow{5}{*}{0.186914} & $-0.1563\ (-1.23\%)$ \\
    & 48 & & $\mathrm{NaN}\ (\mathrm{NaN}\%)$ & & $\mathrm{NaN}\ (\mathrm{NaN}\%)$ & & $\mathrm{NaN}\ (\mathrm{NaN}\%)$ & & $-0.05685\ (-0.46\%)$ & & $-0.1489\ (-1.20\%)$ \\
    & 49 & & $\mathrm{NaN}\ (\mathrm{NaN}\%)$ & & $\mathrm{NaN}\ (\mathrm{NaN}\%)$ & & $\mathrm{NaN}\ (\mathrm{NaN}\%)$ & & $-0.05316\ (-0.44\%)$ & & $-0.1419\ (-1.18\%)$ \\
    & 50 & & $\mathrm{NaN}\ (\mathrm{NaN}\%)$ & & $\mathrm{NaN}\ (\mathrm{NaN}\%)$ & & $\mathrm{NaN}\ (\mathrm{NaN}\%)$ & & $-0.04984\ (-0.42\%)$ & & $-0.1354\ (-1.15\%)$ \\
    & 51 & & $\mathrm{NaN}\ (\mathrm{NaN}\%)$ & & $\mathrm{NaN}\ (\mathrm{NaN}\%)$ & & $\mathrm{NaN}\ (\mathrm{NaN}\%)$ & & $-0.04649\ (-0.41\%)$ & & $-0.1291\ (-1.13\%)$ \\
    \midrule
    \multirow{3}{*}{Case VII} & 95 & \multirow{3}{*}{70.4046528} & $\mathrm{NaN}\ (\mathrm{NaN}\%)$ & \multirow{3}{*}{6.2345658} & $\mathrm{NaN}\ (\mathrm{NaN}\%)$ & \multirow{3}{*}{0.032753} & $\mathrm{NaN}\ (\mathrm{NaN}\%)$ & \multirow{3}{*}{0.2169977} & $-11.64\ (-99.30\%)$ & \multirow{3}{*}{0.235242} & $-11.67\ (-99.56\%)$ \\
    & 100 & & $\mathrm{NaN}\ (\mathrm{NaN}\%)$ & & $\mathrm{NaN}\ (\mathrm{NaN}\%)$ & & $\mathrm{NaN}\ (\mathrm{NaN}\%)$ & & $-8.959\ (-99.57\%)$ & & $-8.974\ (-99.73\%)$ \\
    & 105 & & $\mathrm{NaN}\ (\mathrm{NaN}\%)$ & & $\mathrm{NaN}\ (\mathrm{NaN}\%)$ & & $\mathrm{NaN}\ (\mathrm{NaN}\%)$ & & $-6.692\ (-99.74\%)$ & & $-6.698\ (-99.84\%)$ \\
    \bottomrule
  \end{tabular}%
  }
\end{table}

从表中结果可以看到，精确模拟” 的三类终端方法，仅能在 Case I 基准算例的单一执行价场景下输出有效结果，在其余测试场景中均完全数值失效，不具备全场景工程实用性。

具体而言，Baldeaux (2012) 精确模拟方法仅在 Case I 的单一执行价场景下输出有效定价结果，相对偏差趋近于 0，具备理论层面的无偏性；但在 Case II 至 Case VII 的全部 6 组测试案例、22 个执行价场景中，均出现数值溢出与计算失效，输出结果为非数值型。该现象的核心成因在于，该方法需要逐路径计算高维修正贝塞尔函数与逆拉普拉斯变换，在非基准参数下，极易出现贝塞尔函数数值溢出、积分收敛失败的问题，仅能在文献给定的基准参数下完成计算，无法适配实际市场中的非基准参数环境。基于 Fourier-Cosine 的终端精确模拟方法与基于矩匹配的终端精确模拟方法，表现与 Baldeaux 方法高度一致，仅在 Case I 单一执行价场景下输出有效结果，且相对偏差分别达 $-11.23\%$ 与 $-11.77\%$，定价精度极差，其余所有场景均完全数值失效。两类方法失效的核心成因分别为：Fourier-Cosine 方法对积分方差分布的余弦级数截断，在非基准参数下，分布尾部的截断误差被无限放大，导致特征函数计算失效；矩匹配方法仅匹配分布前四阶矩的 Gram-Charlier 展开，无法适配 3/2 模型方差分布的非正态厚尾特征，易出现负密度的非物理结果，最终导致累积分布函数逆变换抽样失败。

两类近似精确模拟方法中，仅在 Case VI 平值、中等期限、正相关场景下，定价偏差处于可控范围，其余绝大多数场景均出现系统性的大幅定价偏差。在 Case VI 的 5 个近平值执行价场景中，IG 方法的相对偏差处于 $-0.41\%$ 至 $-0.48\%$ 区间，Gamma 方法的相对偏差处于 $-1.13\%$ 至 $-1.23\%$ 区间，偏差可控且计算时间仅为 0.19 秒左右，具备较高的计算效率。但在其余所有测试场景中，两类方法均出现显著的定价偏差：在高 vol-of-vol、强负相关的 Case II、III、VII 场景中，两类方法在所有执行价下的相对偏差幅度达 $-80\%$ 至 $-99.8\%$，定价结果完全偏离基准值；在短期到期、深度实值的 Case IV、V 场景中，Gamma 方法出现 $500\%$ 以上的正向系统性偏差，IG 方法也出现 $-78\%$ 至 $-87\%$ 的负向系统性偏差，完全丧失定价参考价值。该现象的核心成因在于，强负相关、高 vol-of-vol 与短期到期场景下，积分方差的分布呈现极端厚尾与高度右偏特征，仅匹配前两阶矩的双参数 IG 与 Gamma 分布，无法拟合分布的高阶特征与尾部形态，最终导致系统性的定价偏差。

\section{傅里叶频域定价方法结果分析}

频域定价方法完全规避了蒙特卡洛模拟的路径生成过程，通过特征函数的傅里叶积分直接求解期权价格，包含快速傅里叶变换（FFT）与傅里叶余弦（COS）两种方法，完整结果如下表所示：

\begin{table}[htbp]
  \centering
  \caption{Fourier Frequency Domain Pricing Results}
  \label{tab:frequency}
  \resizebox{\textwidth}{!}{%
  \begin{tabular}{lclcccc}
    \toprule
    \multirow{2}{*}{Case} & \multirow{2}{*}{Strike} & \multicolumn{2}{c}{Fast Fourier Transformation} & \multicolumn{2}{c}{Fourier-Cosine Method} \\
    \cmidrule(lr){3-4} \cmidrule(lr){5-6}
    & & Time (s) & Option Bias & Time (s) & Option Bias \\
    \midrule
    Case I & 1 & 0.3843623 & $-2.434\times10^{-7}\ (-0.00\%)$ & 0.1115441 & $-2.437\times10^{-7}\ (-0.00\%)$ \\
    \midrule
    \multirow{3}{*}{Case II}
    & 95 & \multirow{3}{*}{0.1986538} & $-0.0004891\ (-0.00\%)$ & \multirow{3}{*}{0.0517513} & $-0.0004891\ (-0.00\%)$ \\
    & 100 & & $-0.0001461\ (-0.00\%)$ & & $-0.0001461\ (-0.00\%)$ \\
    & 105 & & $-0.0009548\ (-0.02\%)$ & & $-0.0009548\ (-0.02\%)$ \\
    \midrule
    \multirow{3}{*}{Case III}
    & 95 & \multirow{3}{*}{0.2813563} & $-0.0005218\ (-0.00\%)$ & \multirow{3}{*}{0.0735898} & $-0.0005218\ (-0.00\%)$ \\
    & 100 & & $-0.001246\ (-0.01\%)$ & & $-0.001246\ (-0.01\%)$ \\
    & 105 & & $-0.0008511\ (-0.01\%)$ & & $-0.0008511\ (-0.01\%)$ \\
    \midrule
    \multirow{5}{*}{Case IV}
    & 47 & \multirow{5}{*}{15.8604601} & $0.0003048\ (0.00\%)$ & \multirow{5}{*}{3.7037431} & $0.0003048\ (0.00\%)$ \\
    & 48 & & $0.068\ (1.05\%)$ & & $0.068\ (1.05\%)$ \\
    & 49 & & $0.001088\ (0.02\%)$ & & $0.001088\ (0.02\%)$ \\
    & 50 & & $0.001157\ (0.02\%)$ & & $0.001157\ (0.02\%)$ \\
    & 51 & & $0.001138\ (0.02\%)$ & & $0.001138\ (0.02\%)$ \\
    \midrule
    \multirow{3}{*}{Case V}
    & 10 & \multirow{3}{*}{16.0524815} & $0.0015\ (0.00\%)$ & \multirow{3}{*}{3.788367} & $0.0015\ (0.00\%)$ \\
    & 20 & & $0.001557\ (0.01\%)$ & & $0.001557\ (0.01\%)$ \\
    & 40 & & $0.002021\ (0.02\%)$ & & $0.002021\ (0.02\%)$ \\
    \midrule
    \multirow{5}{*}{Case VI}
    & 47 & \multirow{5}{*}{0.2070812} & $0.01801\ (0.14\%)$ & \multirow{5}{*}{0.0686066} & $0.01785\ (0.14\%)$ \\
    & 48 & & $0.01801\ (0.15\%)$ & & $0.01791\ (0.15\%)$ \\
    & 49 & & $0.01807\ (0.15\%)$ & & $0.01791\ (0.15\%)$ \\
    & 50 & & $0.01794\ (0.15\%)$ & & $0.01778\ (0.15\%)$ \\
    & 51 & & $0.01794\ (0.16\%)$ & & $0.01775\ (0.16\%)$ \\
    \midrule
    \multirow{3}{*}{Case VII}
    & 95 & \multirow{3}{*}{0.3207817} & $-2.18\times10^{-5}\ (-0.00\%)$ & \multirow{3}{*}{0.0993632} & $-2.179\times10^{-5}\ (-0.00\%)$ \\
    & 100 & & $-4.598\times10^{-5}\ (-0.00\%)$ & & $-4.597\times10^{-5}\ (-0.00\%)$ \\
    & 105 & & $4.895\times10^{-5}\ (0.00\%)$ & & $4.895\times10^{-5}\ (0.00\%)$ \\
    \bottomrule
  \end{tabular}%
  }
\end{table}

从表中结果可以看到，两类方法是本次实验中唯一在全部 7 组测试案例、所有执行价场景下均无数值失效，同时保持极高定价精度的方法，综合性能显著优于所有蒙特卡洛类方法。

在数值稳定性层面，FFT 与 COS 方法在所有测试场景中均未出现非数值型或无穷大结果，无论是基准平值场景、极端高 vol-of-vol 强负相关场景、短期到期场景，还是深度实值/虚值场景，均可稳定输出定价结果，完全规避了蒙特卡洛类方法的抽样失效与分布拟合失效问题。在定价精度层面，两类方法在所有场景下的相对偏差均控制在 $1.05\%$ 以内，绝大多数场景的相对偏差小于 $0.02\%$。在 Case II 极端高 vol-of-vol 场景中，两类方法的最大相对偏差仅为 $-0.02\%$；在 Case IV 短期到期场景中，最大相对偏差为 $1.05\%$，其余执行价的偏差均控制在 $0.02\%$ 以内；在 Case VI 高 vol-of-vol 场景中，相对偏差仅为 $0.15\%$ 左右，定价精度远超所有蒙特卡洛类方法。该结果的核心成因在于，频域方法基于 3/2 模型特征函数的闭式解，通过数值积分直接完成定价，不存在蒙特卡洛方法的抽样误差、分布拟合误差，也不存在时间离散的截断误差，仅存在可忽略的数值积分截断误差。

在计算效率层面，COS 方法的计算效率全面领先于 FFT 方法，在所有测试场景中，COS 方法的运行时间仅为 FFT 方法的 $1/3$ 至 $1/4$。在 Case IV 多执行价场景中，FFT 方法的运行时间为 15.86 秒，COS 方法仅为 3.70 秒；在单执行价场景中，COS 方法的运行时间仅为 0.05 至 0.11 秒，计算效率远高于所有蒙特卡洛类方法。该效率差异的核心理论成因在于，COS 方法基于余弦级数展开具备指数级收敛速度，仅需 128 项截断即可达到 FFT 方法 4096 点网格的计算精度，同时不存在 FFT 方法的栅栏效应，可灵活计算任意执行价的期权价格，无需额外的插值计算，进一步降低了计算成本。

\section{综合性能对比与研究结论}

基于 7 组测试案例的全场景实验结果，可从数值稳定性、定价精度、计算效率三个核心维度，对九类定价方法的综合性能进行系统性分层。第一梯队为 COS 与 FFT 频域方法，是唯一具备全场景数值稳定性、同时保持极高定价精度与计算效率的方法，其中 COS 方法的综合性能全面领先 FFT 方法，为 3/2 模型欧式期权定价的最优工程方案。第二梯队为时间步进类的精确步进、Poisson-Gamma 近似精确步进与 QE 近似精确步进方法，是蒙特卡洛框架下唯一具备全场景数值稳定性的方案，可稳定输出完整的方差路径，定价精度与计算效率平衡良好，为路径依赖期权定价、动态对冲等需要完整路径生成场景的唯一可靠方案。第三梯队为 Euler/Milstein 离散格式，仅在细步长、常规参数场景下有效，极端场景下完全数值失效，计算效率极低，仅可作为理论精度对比的基准方案。第四梯队为 IG 与 Gamma 近似精确模拟方法，仅在特定平值场景下偏差可控，绝大多数场景出现系统性大幅定价偏差，仅可用于特定场景的学术对比，不具备工程实用性。第五梯队为 Baldeaux 精确模拟、Fourier-Cosine 终端模拟与矩匹配终端模拟方法，仅能在单一基准场景下运行，其余所有场景均完全数值失效，仅可用于理论层面的无偏性验证。

基于全场景的实验验证，可得到三个结论。第一，3/2 模型定价方法的理论无偏性，不等同于工程应用中的全场景稳定性，标称 “精确” 的终端模拟方法，在非基准参数下极易出现特殊函数数值溢出、积分收敛失败等问题，仅能在文献基准场景下完成理论验证，无法适配实际市场的复杂参数环境。第二，基于方差过程分布特性构造的时间步进抽样方法，是蒙特卡洛框架下的最优实现方案，此类方法天然保证了方差过程的非负性，在所有极端场景下均保持数值稳定性，同时定价精度显著优于传统 Euler 离散格式，为需要路径生成的衍生品定价场景提供了可靠的实现路径。第三，基于特征函数的 COS 频域方法，是 3/2 模型欧式期权定价的绝对最优方案，该方法在所有测试场景下均保持了极高的精度与稳定性，无任何数值失效，同时计算效率远超所有蒙特卡洛类方法，在欧式期权批量定价、模型参数校准、极端市场环境定价等场景中，具备绝对的性能优势。

基于全场景的性能验证，我们给出不同的应用场景对应方法推荐。对于欧式期权批量定价、模型参数校准场景，优先选择 COS 频域方法，其全场景稳定性、最高精度与最优效率可充分适配此类场景的需求；对于路径依赖期权定价、动态对冲等需要完整方差路径的场景，优先选择 Poisson-Gamma 或 QE 近似精确步进方法，其全场景稳定性与精度效率平衡可满足此类场景的工程需求；对于理论精度基准对比场景，仅可在基准参数环境下使用 Baldeaux 精确模拟方法；其余终端类方法不推荐在实际定价场景中使用，仅可用于特定场景的学术对比分析。